\begin{center}\begin{tikzpicture}[font=\small,>=stealth,line join=round]
\draw[gray!40,rounded corners=3pt] (0,0) rectangle (14.2,8.3);\node at (7.1,7.700000000000001) {分式证明一：输出再作为输入，回到原来的数};
\begin{scope}[shift={(2.975,3.7409090909090907)},x=0.75cm,y=0.41818181818181815cm]\draw[->,gray] (-2.5,0)--(5.5,0) node[right]{$u$};\draw[->,gray] (0,-5)--(0,6) node[above]{$v$};\draw[blue,thick,solid,domain=-2:0.6,samples=180] plot(\x,{(\x+1)/(\x-1)});\draw[blue,thick,solid,domain=1.5:5,samples=180] plot(\x,{(\x+1)/(\x-1)});\draw[gray,thick,dashed,domain=-2:5,samples=180] plot(\x,{\x});\fill[orange!85!black] (2,3) circle(2.2pt);\fill[orange!85!black] (3,2) circle(2.2pt);\draw[orange!85!black,->] (2,3)--(3,3)--(3,2)--(2,2);\end{scope}\node[anchor=west,align=left,text width=5.0cm] at (8.4,4.0) {$\Phi(u)=\frac{u+1}{u-1}$\\[9pt]图中 $2\to3\to2$\\[9pt]蓝曲线为 $v=\Phi(u)$\\[9pt]虚线为 $v=u$\\[9pt]一般需证 $\Phi(\Phi(u))=u$};
\node[text width=12.6cm,align=center] at (7.1,.65) {橙线把第一次输出移作第二次输入；字母通分证明覆盖所有 $u\ne1$。};
\end{tikzpicture}\end{center}
